% ============================================================
\chapter*{Preface}
\addcontentsline{toc}{chapter}{Preface}
\markboth{Preface}{Preface}
% ============================================================

% ------------------------------------------------------------
\section*{Course Objectives}
% ------------------------------------------------------------

\emph{Natural Language Processing} (NLP) is one of the most dynamic areas of
contemporary artificial intelligence. This graduate-level course (Master/PhD)
covers classical mathematical foundations
through to recent large language models.

The course has three primary objectives:
\begin{enumerate}
    \item \textbf{Understand theoretical foundations} --- probabilistic language
          modeling, information theory, linear algebra for vector
          representations.
    \item \textbf{Master modern architectures} --- recurrent networks, attention
          mechanisms, Transformers, pre-trained models (BERT, GPT, T5).
    \item \textbf{Develop critical thinking} --- rigorous evaluation, ethical
          considerations, algorithmic bias, generative system safety.
\end{enumerate}

% ------------------------------------------------------------
\section*{Target Audience}
% ------------------------------------------------------------

This course is designed for Master's and PhD students in computer science,
data science, or computational linguistics. The following prerequisites are
assumed:

\begin{itemize}
    \item \textbf{Mathematics}: linear algebra (vector spaces, matrix
          decompositions), multivariate calculus, probability and statistics.
    \item \textbf{Computer science}: Python programming, basic algorithms,
          familiarity with scientific libraries (NumPy, PyTorch or TensorFlow).
    \item \textbf{Machine learning}: regression, classification, gradient
          descent, basic neural networks.
\end{itemize}

% ------------------------------------------------------------
\section*{Course Organization}
% ------------------------------------------------------------

The course is organized into eleven chapters, following a logical progression
from classical text representations to the most recent generative systems:

\begin{description}
    \item[Chapter~1 --- Introduction to NLP.]
        Historical overview, fundamental tasks, challenges specific to natural
        language.
    \item[Chapter~2 --- Classical Text Representations.]
        Bag-of-words models, TF-IDF, $n$-grams, sparse representations.
    \item[Chapter~3 --- Word Embeddings.]
        Word2Vec (CBOW, Skip-gram), GloVe, FastText; geometric properties of
        embedding spaces.
    \item[Chapter~4 --- Sequence Models.]
        Recurrent neural networks (RNN), LSTM, GRU; vanishing and exploding
        gradient problems.
    \item[Chapter~5 --- Attention and Transformers.]
        Attention mechanism, self-attention, complete Transformer architecture.
    \item[Chapter~6 --- Pre-trained Models.]
        BERT, GPT, T5; the pre-training/fine-tuning paradigm.
    \item[Chapter~7 --- Fine-tuning and Instruction Tuning.]
        Adaptation strategies, LoRA, RLHF, alignment.
    \item[Chapter~8 --- Text Generation and Decoding.]
        Greedy decoding, beam search, sampling, top-$k$, top-$p$.
    \item[Chapter~9 --- Retrieval-Augmented Generation (RAG).]
        Passage retrieval, external knowledge integration, RAG pipelines.
    \item[Chapter~10 --- LLM Evaluation.]
        Automatic metrics, human evaluation, benchmarks, robustness.
    \item[Chapter~11 --- Ethics, Bias, and Safety.]
        Data and model bias, toxicity, privacy, regulation.
\end{description}

% ------------------------------------------------------------
\section*{Typographic Conventions}
% ------------------------------------------------------------

\begin{itemize}
    \item Formal \textbf{definitions} are framed in \texttt{definition}
          environments.
    \item \textbf{Theorems}, \textbf{propositions}, and \textbf{lemmas} are
          numbered and, when relevant, accompanied by a proof.
    \item \textbf{Intuition} boxes provide qualitative understanding before
          formalization.
    \item \textbf{Warning} boxes highlight common mistakes.
    \item \textbf{Best practice} boxes summarize practical recommendations.
    \item Python code primarily uses the \texttt{transformers} (Hugging Face),
          \texttt{torch}, and \texttt{scikit-learn} libraries.
\end{itemize}

% ------------------------------------------------------------
\section*{Mathematical Notation}
% ------------------------------------------------------------

We adopt the following notation throughout the course:

\begin{center}
\renewcommand{\arraystretch}{1.3}
\begin{tabular}{cl}
\toprule
\textbf{Notation} & \textbf{Meaning} \\
\midrule
$\R^d$ & Euclidean space of dimension $d$ \\
$\N$ & Set of natural numbers \\
$\mathbf{x}, \mathbf{W}$ & Vectors (bold lowercase), matrices (bold uppercase) \\
$\inner{\mathbf{u}}{\mathbf{v}}$ & Inner product of $\mathbf{u}$ and $\mathbf{v}$ \\
$\norm{\mathbf{x}}$ & Euclidean norm of $\mathbf{x}$ \\
$\abs{S}$ & Cardinality of set $S$ \\
$\E[X]$ & Expectation of random variable $X$ \\
$P(A \mid B)$ & Conditional probability of $A$ given $B$ \\
$\KL(p \| q)$ & Kullback-Leibler divergence from $p$ to $q$ \\
$\sigma(\cdot)$ & Logistic sigmoid function \\
$\softmax(\cdot)$ & Softmax function \\
$\argmax_x f(x)$ & Argument maximizing $f$ \\
$\argmin_x f(x)$ & Argument minimizing $f$ \\
$\mathcal{V}$ & Vocabulary (finite set of tokens) \\
$V = \abs{\mathcal{V}}$ & Vocabulary size \\
$\mathbf{E} \in \R^{V \times d}$ & Embedding matrix \\
$T$ & Sequence length \\
\bottomrule
\end{tabular}
\end{center}

% ------------------------------------------------------------
\section*{Software and Environment}
% ------------------------------------------------------------

Labs and code examples rely on the following ecosystem:

\begin{itemize}
    \item \textbf{Python} $\geq 3.10$
    \item \textbf{PyTorch} $\geq 2.0$
    \item \textbf{Hugging Face Transformers} $\geq 4.35$
    \item \textbf{Hugging Face Datasets}
    \item \textbf{scikit-learn} for classical methods
    \item \textbf{NLTK} and \textbf{spaCy} for linguistic preprocessing
    \item \textbf{Matplotlib} / \textbf{seaborn} for visualization
\end{itemize}

\begin{bestpractice}[title=\textbf{Reproducible environment}]
We recommend using a virtual environment (\texttt{venv} or \texttt{conda}) and
fixing random seeds (\texttt{torch.manual\_seed}, \texttt{random.seed},
\texttt{numpy.random.seed}) to ensure experiment reproducibility.
\end{bestpractice}

% ------------------------------------------------------------
\section*{Historical Perspective}
% ------------------------------------------------------------

NLP has undergone several paradigm shifts. In the 1950s, early approaches were
based on manually coded linguistic rules. The 1990s saw the rise of statistical
methods (hidden Markov models, $n$-gram models). The 2010s were dominated by
deep learning (word embeddings, recurrent networks). Since 2017, the Transformer
architecture has ushered in the era of large language models (LLMs), radically
transforming the landscape of the field.

\begin{intuition}[title=\textbf{Why is language hard?}]
Natural language is inherently ambiguous (polysemy, anaphora, ellipsis),
contextual (meaning depends on pragmatic context), and compositional (sentence
meaning depends on syntactic structure and constituent meanings). These
properties make NLP fundamentally different from processing structured data.
\end{intuition}

% ------------------------------------------------------------
\section*{How to Use This Course}
% ------------------------------------------------------------

Each chapter is designed to be self-contained while building on concepts from
previous chapters. We recommend sequential reading for a first pass, then
targeted use as a reference.

Exercises at the end of each chapter are classified into three difficulty levels:
\begin{itemize}
    \item[$\star$] Comprehension and verification exercises.
    \item[$\star\star$] Deepening exercises requiring mathematical reasoning or
          implementation.
    \item[$\star\star\star$] Open research exercises that can serve as starting
          points for projects or theses.
\end{itemize}

% ------------------------------------------------------------
\section*{Acknowledgments}
% ------------------------------------------------------------

This course has benefited from the contributions of many colleagues,
researchers, and students. We particularly thank the open-source communities
of Hugging Face, PyTorch, and spaCy, whose tools make NLP teaching both
rigorous and accessible.

We also thank students from previous cohorts whose questions and feedback have
greatly improved the quality of this pedagogical material.

\vfill
\begin{flushright}
\textit{[Author Name], March 2026}
\end{flushright}
